\documentclass[a4paper,11pt,oneside]{memoir}

% Castellano
\usepackage[spanish,es-tabla]{babel}
\selectlanguage{spanish}
\usepackage[utf8]{inputenc}
\usepackage{placeins}

\RequirePackage{booktabs}
\RequirePackage[table]{xcolor}
\RequirePackage{xtab}
\RequirePackage{multirow}

% Links
\usepackage[colorlinks]{hyperref}
\hypersetup{
	allcolors = {blue}
}

% Ecuaciones
\usepackage{amsmath}

% Rutas de fichero / paquete
\newcommand{\ruta}[1]{{\sffamily #1}}

% Párrafos
\nonzeroparskip

% Símbolo euro
\usepackage{eurosym}

% Imagenes
\usepackage{graphicx}
\newcommand{\imagen}[2]{
	\begin{figure}[!h]
		\centering
		\includegraphics[width=0.9\textwidth]{#1}
		\caption{#2}\label{fig:#1}
	\end{figure}
	\FloatBarrier
}

\newcommand{\imagenflotante}[2]{
	\begin{figure}%[!h]
		\centering
		\includegraphics[width=0.9\textwidth]{#1}
		\caption{#2}\label{fig:#1}
	\end{figure}
}

\newcommand{\imagenLarga}[2]{
	\begin{figure}[!h]
		\centering
		\includegraphics[width=1.2\textwidth]{#1}
		\caption{#2}\label{fig:#1}
	\end{figure}
	\FloatBarrier
}

\newcommand{\imagenMediana}[2]{
	\begin{figure}[!h]
		\centering
		\includegraphics[width=0.6\textwidth]{#1}
		\caption{#2}\label{fig:#1}
	\end{figure}
	\FloatBarrier
}

\newcommand{\imagenPeque}[2]{
	\begin{figure}[!h]
		\centering
		\includegraphics[width=0.4\textwidth]{#1}
		\caption{#2}\label{fig:#1}
	\end{figure}
	\FloatBarrier
}

\newcommand{\imagenMuyPeque}[2]{
	\begin{figure}[!h]
		\centering
		\includegraphics[width=0.2\textwidth]{#1}
		\caption{#2}\label{fig:#1}
	\end{figure}
	\FloatBarrier
}

\usepackage{color}
\usepackage{colortbl}
\usepackage{multicol}
\usepackage{multirow}
\usepackage{booktabs}
\usepackage{tabularx}
\usepackage{array}
\usepackage{float}



% El comando \figura nos permite insertar figuras comodamente, y utilizando
% siempre el mismo formato. Los parametros son:
% 1 -> Porcentaje del ancho de página que ocupará la figura (de 0 a 1)
% 2 --> Fichero de la imagen
% 3 --> Texto a pie de imagen
% 4 --> Etiqueta (label) para referencias
% 5 --> Opciones que queramos pasarle al \includegraphics
% 6 --> Opciones de posicionamiento a pasarle a \begin{figure}
\newcommand{\figuraConPosicion}[6]{%
  \setlength{\anchoFloat}{#1\textwidth}%
  \addtolength{\anchoFloat}{-4\fboxsep}%
  \setlength{\anchoFigura}{\anchoFloat}%
  \begin{figure}[#6]
    \begin{center}%
      \Ovalbox{%
        \begin{minipage}{\anchoFloat}%
          \begin{center}%
            \includegraphics[width=\anchoFigura,#5]{#2}%
            \caption{#3}%
            \label{#4}%
          \end{center}%
        \end{minipage}
      }%
    \end{center}%
  \end{figure}%
}

%
% Comando para incluir imágenes en formato apaisado (sin marco).
\newcommand{\figuraApaisadaSinMarco}[5]{%
  \begin{figure}%
    \begin{center}%
    \includegraphics[angle=90,height=#1\textheight,#5]{#2}%
    \caption{#3}%
    \label{#4}%
    \end{center}%
  \end{figure}%
}
% Para las tablas
\newcommand{\otoprule}{\midrule [\heavyrulewidth]}
%
% Nuevo comando para tablas pequeñas (menos de una página).
\newcommand{\tablaSmall}[5]{%
 \begin{table}[h!]
  \begin{center}
   \rowcolors {2}{gray!35}{}
   \begin{tabular}{#2}
    \toprule
    #4
    \otoprule
    #5
    \bottomrule
   \end{tabular}
   \caption{#1}
   \label{tabla:#3}
  \end{center}
 \end{table}
}

%
% Nuevo comando para tablas pequeñas (menos de una página).
\newcommand{\tablaSmallSinColores}[5]{%
 \begin{table}[H]
  \begin{center}
   \begin{tabular}{#2}
    \toprule
    #4
    \otoprule
    #5
    \bottomrule
   \end{tabular}
   \caption{#1}
   \label{tabla:#3}
  \end{center}
 \end{table}
}

\newcommand{\tablaApaisadaSmall}[5]{%
\begin{landscape}
  \begin{table}
   \begin{center}
    \rowcolors {2}{gray!35}{}
    \begin{tabular}{#2}
     \toprule
     #4
     \otoprule
     #5
     \bottomrule
    \end{tabular}
    \caption{#1}
    \label{tabla:#3}
   \end{center}
  \end{table}
\end{landscape}
}

%
% Nuevo comando para tablas grandes con cabecera y filas alternas coloreadas en gris.
\newcommand{\tabla}[6]{%
  \begin{center}
    \tablefirsthead{
      \toprule
      #5
      \otoprule
    }
    \tablehead{
      \multicolumn{#3}{l}{\small\sl continúa desde la página anterior}\\
      \toprule
      #5
      \otoprule
    }
    \tabletail{
      \hline
      \multicolumn{#3}{r}{\small\sl continúa en la página siguiente}\\
    }
    \tablelasttail{
      \hline
    }
    \bottomcaption{#1}
    \rowcolors {2}{gray!35}{}
    \begin{xtabular}{#2}
      #6
      \bottomrule
    \end{xtabular}
    \label{tabla:#4}
  \end{center}
}

%
% Nuevo comando para tablas grandes con cabecera.
\newcommand{\tablaSinColores}[6]{%
	\begin{center}
		\tablefirsthead{
			\toprule
			#5
			\otoprule
		}
		\tablehead{
			\multicolumn{#3}{l}{\small continúa desde la página anterior}\\
			\toprule
			#5
			\otoprule
		}
		\tabletail{
			\hline
			\multicolumn{#3}{r}{\small continúa en la página siguiente}\\
		}
		\tablelasttail{
			\hline
		}
		\bottomcaption{#1}
		\begin{xtabular}{#2}
			#6
			\bottomrule
		\end{xtabular}
		\label{tabla:#4}
	\end{center}
}

%
% Nuevo comando para tablas grandes sin cabecera.
\newcommand{\tablaSinCabecera}[5]{%
  \begin{center}
    \tablefirsthead{
      \toprule
    }
    \tablehead{
      \multicolumn{#3}{l}{\small\sl continúa desde la página anterior}\\
      \hline
    }
    \tabletail{
      \hline
      \multicolumn{#3}{r}{\small\sl continúa en la página siguiente}\\
    }
    \tablelasttail{
      \hline
    }
    \bottomcaption{#1}
  \begin{xtabular}{#2}
    #5
   \bottomrule
  \end{xtabular}
  \label{tabla:#4}
  \end{center}
}



\definecolor{cgoLight}{HTML}{EEEEEE}
\definecolor{cgoExtralight}{HTML}{FFFFFF}

%
% Nuevo comando para tablas grandes sin cabecera.
\newcommand{\tablaSinCabeceraConBandas}[5]{%
  \begin{center}
    \tablefirsthead{
      \toprule
    }
    \tablehead{
      \multicolumn{#3}{l}{\small\sl continúa desde la página anterior}\\
      \hline
    }
    \tabletail{
      \hline
      \multicolumn{#3}{r}{\small\sl continúa en la página siguiente}\\
    }
    \tablelasttail{
      \hline
    }
    \bottomcaption{#1}
    \rowcolors[]{1}{cgoExtralight}{cgoLight}

  \begin{xtabular}{#2}
    #5
   \bottomrule
  \end{xtabular}
  \label{tabla:#4}
  \end{center}
}




\graphicspath{ {./img/} }

% Capítulos
\chapterstyle{bianchi}
\newcommand{\capitulo}[2]{
	\setcounter{chapter}{#1}
	\setcounter{section}{0}
	\chapter*{#2}
	\addcontentsline{toc}{chapter}{#2}
	\markboth{#2}{#2}
}

% Apéndices
\renewcommand{\appendixname}{Apéndice}
\renewcommand*\cftappendixname{\appendixname}

\newcommand{\apendice}[1]{
	%\renewcommand{\thechapter}{A}
	\chapter{#1}
}

\renewcommand*\cftappendixname{\appendixname\ }

% Formato de portada
\makeatletter
\usepackage{xcolor}
\newcommand{\tutor}[1]{\def\@tutor{#1}}
\newcommand{\course}[1]{\def\@course{#1}}
\definecolor{cpardoBox}{HTML}{E6E6FF}
\def\maketitle{
  \null
  \thispagestyle{empty}
  % Cabecera ----------------
\noindent\includegraphics[width=\textwidth]{cabecera}\vspace{1cm}%
  \vfill
  % Título proyecto y escudo informática ----------------
  \colorbox{cpardoBox}{%
    \begin{minipage}{.8\textwidth}
      \vspace{.5cm}\Large
      \begin{center}
      \textbf{TFG del Grado en Ingeniería Informática}\vspace{.6cm}\\
      \textbf{\LARGE\@title{}}
      \end{center}
      \vspace{.2cm}
    \end{minipage}

  }%
  \hfill\begin{minipage}{.20\textwidth}
    \includegraphics[width=\textwidth]{escudoInfor}
  \end{minipage}
  \vfill
  % Datos de alumno, curso y tutores ------------------
  \begin{center}%
  {%
    \noindent\LARGE
    Presentado por \@author{}\\ 
    en Universidad de Burgos --- \@date{}\\
    Tutor: \@tutor{}\\
  }%
  \end{center}%
  \null
  \cleardoublepage
  }
\makeatother


% Datos de portada
\title{UBUassistant \\Documentación Técnica}
\author{Carlos González Calatrava}
\tutor{Pedro Renedo Fernández}
\date{\today}

\begin{document}

\maketitle



\cleardoublepage



%%%%%%%%%%%%%%%%%%%%%%%%%%%%%%%%%%%%%%%%%%%%%%%%%%%%%%%%%%%%%%%%%%%%%%%%%%%%%%%%%%%%%%%%



\frontmatter


\clearpage

% Indices
\tableofcontents

\clearpage

\listoffigures

\clearpage

\listoftables

\clearpage

\mainmatter

\appendix

\apendice{Plan de Proyecto Software}

\section{Introducción}
Una correcta planificación del proyecto es muy importante. Esto se debe por que de ello dependerá de que nuestro proyecto vaya evolucionando de una forma correcta y que se cumplan los tiempos de entrega, evitando así todos los retrasos posibles. Esta fase de planificación está compuesta por dos apartados:

\begin{itemize}
	\tightlist
	\item
	\textbf{Planificación temporal:} esta planificación consta principalmente en establecer los diferentes tiempos del proyecto. En este apartado se establece la fecha de inicio y la fecha de finalización del proyecto. A parte se establece la duración de cada una de las fases del proyecto, aunque estas pueden adaptarse algo a lo largo del proyecto en función del trabajo realizado a lo largo de cada uno de los sprints.
	
	\item
	\textbf{Estudio de viabilidad:} este estudio es muy importante, ya que es lo que nos va a permitir valorar si el proyecto nos va a permitir continuar hacia delante con él o no. Dentro de este estudio nos encontramos con dos fases;
	\begin{itemize}
		\tightlist
		\item
		\textbf{Viabilidad económica:} es el estudio en el que vamos a estudiar los costes del proyecto, tanto de tiempo, como de personal así como los recursos de terceros que son necesarios para realizarlo.
		\item
		\textbf{Viabilidad legal:} este estudio es el más importante de todos, ya que es el que necesitamos para saber si el proyecto que vamos a hacer está permitido por ley. A parte de esto, también se estudia todo lo relacionado con la adquisición de las diferentes licencias que sean necesarias para el desarrollo del proyecto.
	\end{itemize}
\end{itemize}

\section{Planificación temporal}
Para el desarrollo del proyecto se optó por seguir la metodología de desarrollo ágil \textit{Scrum} \cite{scrum:wiki}.

Debido al tamaño del proyecto realizado, no se ha podido seguir a la totalidad los rasgos de esta metodología, pero se ha intentado adaptar todo lo posible, manteniendo la esencia de la misma. Las diferentes pausas que se siguieron fueron las siguientes:

\begin{itemize}
	\tightlist
	\item
	Desarrollo incremental de la aplicación. Para ello se realizó a través de sprints.
	\item 
	Los sprints se realizaron de una duración de una semana cada uno.
	\item 
	Reuniones al final de cada uno de los sprint, en los cuales se valoraba el trabajo realizado en el sprint y se establecían los objetivos para el próximo sprint.
\end{itemize}

Las tareas realizadas en cada uno de los sprints se comenta a continuación.

\subsection{Sprint 0 (31/02/2018 - 06/03/2018)}

En esta primera reunión lo que se hizo fue hablar sobre la propuesta del tutor para realizar el TFG. En ella se me comentó en que consistió el TFG que me ha tocado mejorar, como un poco las ideas que se quería para mejorar el citado TFG.

La tarea que se me encomendó en este primer sprint fue el de que hiciera una breve investigación sobre el funcionamiento de los algoritmos de razonamiento basado en casos.

\subsection{Sprint 1 (07/03/2018 - 13/03/2018)}

En esta reunión el tutor me ofreció todo el material relacionado con el TFG a mejorar, así como una definición formal de los requisitos a cumplir durante el desarrollo del proyecto.

El objetivo que se estableció en este sprint es el de leerme la documentación del TFG e ir empezando la comprensión del proyecto para ir recogiendo todas las dudas posibles para que pudieran ser respondidas las dudas en la próxima reunión.

\subsection{Sprint 2 (14/03/2018 - 20/03/2018)}

En esta reunión, lo primero que se hizo, fue solucionar todas aquellas dudas que fueron surgiendo a lo largo del sprint anterior y se establecieron los objetivos para el próximo sprint.

El objetivo de este sprint fue el de analizar todas aquellas herramientas que creyera que fueran necesarias para el desarrollo del proyecto y dejar constancia de porque la decisión de emplear esas herramientas en vez de usar otras herramientas.

\subsection{Sprint 3 (21/03/2018 - 27/03/2018)}

Lo primero que se realizó fue valorar las herramientas elegidas para el proyecto, así como alguna recomendación de otras herramientas por parte del tutor, dejando la valoración de esas herramientas para el próximo sprint.

A parte de ese objetivo, también se estableció como objetivo investigar la forma de implementar la introducción de texto mediante voz en nuestra aplicación de Android.

\subsection{Sprint 4 (28/03/2018 - 10/04/2018)}

Al comienzo de la reunión se terminó de valorar las herramientas sugeridas por el profesor en la reunión anterior. Además también se valoró la decisión tomada respecto al reconocimiento de voz dentro de la aplicación, tomando como buena la solución encontrada para ello.

Debido a que este sprint duró dos semanas a raíz de que la Universidad de Burgos cerraba por vacaciones de Semana Santa se establecieron más objetivos para cubrir trabajo en las dos semanas. El primer objetivo que se estableció fue el de investigar la viabilidad de realizar el cambio de framework de razonamiento basado en casos, ya que es uno de los objetivos del proyecto. A parte de esto, también se estableció como objetivo que fuera 
creando el repositorio en GitHub y que se fuera importando el anterior TFG para poder empezar a realizar cambios en el código.






















\section{Estudio de viabilidad}

\subsection{Viabilidad económica}

\subsection{Viabilidad legal}



\apendice{Especificación de Requisitos}

\section{Introducción}

\section{Objetivos generales}

\section{Catalogo de requisitos}

\section{Especificación de requisitos}



\apendice{Especificación de diseño}

\section{Introducción}

En este apartado se va a detallar como se han estructurado los datos, los procedimientos y la arquitectura del proyecto que se ha desarrollado.

\section{Diseño de datos}

En este proyecto, nos podemos encontrar con dos entidades que se relacionan entre ellas.

\begin{itemize}
	\tightlist
	\item
	\textbf{CaseDescription:} esta entidad es la encargada de almacenar los diferentes casos descritos para el algoritmo de IA que lleva el asistente. Para definir el caso lleva un identificador, así como siete palabras clave que describen cada uno de los diferentes casos.
	\item
	\textbf{CaseSolution:} esta entidad se encarga de almacenar las diferentes respuestas que da el algoritmo de IA a cada uno de los diferentes casos descriptivos. Los atributos de esta entidad son un identificador y una respuesta, que por defecto en todos los casos es una dirección web.
\end{itemize}

La relación entre estas dos entidades está realizada a través de Hibernate, sin necesidad de emplear ninguna clave externa para poder hacerlo.

\imagen{DiagramaER.png}{Diagrama E/R}

A parte de las dos entidades que hemos mencionado con anterioridad, existen otras entidades representadas en forma de tablas.

\begin{itemize}
	\tightlist
	\item
	\textbf{aprendizaje:} Esta entidad se encarga de almacenar los diferentes casos que puede aprender nuestro algoritmo. Para almacenar los diferentes casos, se emplea un identificador (\textit{id}), el identificador del usuario que ha propuesto el aprendizaje (\textit{userid}) y dos palabras que definen el caso (\textit{palabra1} \& \textit{palabra2}). 
	\item
	\textbf{logger:} La función de esta entidad es almacenar un registro de todas las respuestas accedidas por los usuarios. Para ello, se utiliza un identificador (\textit{id}), el identificador del usuario que ha accedido a la respuesta (\textit{userid}), la fecha de acceso a la respuesta (\textit{fecha}), las siete posibles palabras que pueden describir el caso (\textit{keyWord1 - keyWord7}), la categoría de la respuesta dada (\textit{categoria}), la respuesta accedida (respuesta), el número de veces que se ha accedido a la respuesta (\textit{num\_busquedas}), el número de veces que se ha votado esa respuesta (\textit{num\_votos}) y la valoración total que ha recibido la respuesta (\textit{valoracion\_total}).
	\item 
	\textbf{frases:} Esta entidad se encarga de almacenar una serie de frases que se otorga junto a la respuesta, para dar así un aspecto más humano a nuestro asistente. Para ello, se utiliza un identificador (\textit{id}) y la frases que se ofrecen junto a la pregunta (\textit{frase}).
	\item 
	\textbf{saludos:}
\end{itemize}

\section{Diseño procedimental}

\section{Diseño arquitectónico}



\apendice{Documentación técnica de programación}

\section{Introducción}

En la documentación técnica del programador se explica la estructura final de directorios de la aplicación desarrollada, así como los pasos para compilar, instalar y ejecutar el proyecto. Además se incluye un manual para el programador y se lleva a cabo una explicación de las pruebas realizadas al sistema.

\section{Estructura de directorios}

El proyecto consta de los siguientes directorios:

\begin{itemize}
\tightlist
\item \textbf{/:} el directorio raíz del proyecto contiene el fichero de la licencia, el archivo README y el fichero que almacena las dependencias que se cargarán dinámicamente mediante \emph{Maven}. Además, se encuentran todas las carpetas que se especifican a continuación.
\item \textbf{/doc:} documentación del proyecto.
\item \textbf{/doc/javadoc:} documentación del código generada por JavaDoc.
\item \textbf{/doc/latex:} documentación generada mediante \emph{LaTeX}.
\item \textbf{/doc/latex/img:} imágenes utilizadas en la documentación.
\item \textbf{/doc/latex/tex:} ficheros utilizados en para generar la documentación \emph{LaTeX}.
\item \textbf{/lib:} archivos binarios y recursos del algoritmo de inteligencia artificial para añadir al \emph{path} del proyecto y hacerlo funcionar.
\item \textbf{/rsc:} recursos del proyecto, como el fichero de creación de base de datos y los distintos ejecutables para realizar las pruebas con \emph{Selenium}.
\item \textbf{/src:} código fuente de la aplicación web.
\item \textbf{/src/main:} clases, recursos y páginas HTML/JSP del proyecto.
\item \textbf{/src/main/java:} clases Java del proyecto y fichero \emph{log4j.properties} que configura el logger.
\item \textbf{/src/main/resources:} recursos de configuración.
\item \textbf{/src/main/webapp:} ficheros para la página web (HTML, JSP, CSS, JS, Imágenes, XML)
\item \textbf{/src/test/java: } test unitarios, de integración y de interfaz.
\item \textbf{/target:} clases compiladas de nuestro proyecto.
\end{itemize}

\section{Manual del programador}

En el proyecto se han utilizado varias herramientas que son necesarias tener instaladas en el sistema para trabajar con el proyecto.

\begin{itemize}
\tightlist
\item \textbf{Java JDK 8.}
\item \textbf{Eclipse JEE.}
\item \textbf{Git.}
\item \textbf{Tomcat.}
\end{itemize}

\subsection{Java JDK 8}

JDK \emph{(Java Development Kit)} son un conjunto de herramientas de desarrollo que permiten crear una aplicación Java. Para su descarga desde la página web oficial \cite{down:jdk} hay que seleccionar la versión deseada y la distribución y arquitectura del sistema operativo adecuada. Una vez realizada la descarga se procede con la instalación siguiendo los pasos del asistente. Finalizada la instalación hay que configurar las variables de entorno. 

Existen tres variables de entorno a configurar \cite{wiki:jdk}

\begin{itemize}
\tightlist
\item \textbf{JAVAPATH:} ruta completa de la instalación de JDK.
\item \textbf{CLASSPATH:} ruta que contiene las bibliotecas o clases del usuario.
\item \textbf{PATH:} variable de entorno donde hay que agregar la ruta de JDK.
\end{itemize}

\subsection{Eclipse JEE}

Eclipse JEE es la distribución del IDE de Eclipse con las herramientas necesarias para el desarrollo web en Java. Este paquete de Eclipse contiene herramientas como Git, Maven, Java EE Developer Tools, JavaScript Developer Tools y demás que hacen posible la programación de una aplicación web \cite{wiki:ecli}.

La descarga se realiza desde la página web oficial \cite{down:ecli} seleccionando la versión Eclipse JEE junto con la arquitectura adecuada del sistema operativo. Una vez descargado se descomprimen los archivos en la ruta deseada.

\imagen{Eclipse.png}{Interfaz Elipse JEE.}

\subsection{Git}

Para comunicarse con el repositorio de control de versiones GitHub, es necesario tener instalado Git. Al utilizar el IDE Eclipse tenemos dos posibilidades de instalación, instalar Git de manera independiente y utilizar comando en consola para trabajar, o optar por instalar el paquete Git que se integra en Eclipse y trabajar mediante una interfaz gráfica.

\imagen{GitBash.png}{Terminal Git Bash.}

\imagen{GitEclipse.png}{Interfaz de Git en Eclipse.}

La descarga de la versión independiente se puede realizar desde la página oficial \cite{down:git}.

La descarga de la versión integrada en Eclipse puede realizarse desde el \emph{MarketPlace}. Se recomienda instalar EGit. Las distribuciones de Eclipse suelen traer instalado de forma automática Git.

\imagen{GitEclipseInstallation}{Descargar de Git desde Eclipse MarketPlace.}

\subsection{Tomcat}

Para realizar las pruebas y poder ejecutar el proyecto durante su desarrollo es necesario disponer de un servidor que incorpore una máquina virtual de Java. Tomcat es un contenedor de aplicaciones Java que puede actuar como servidor para aplicaciones web \cite{wiki:tomc}. Su descarga puede realizarse desde la página oficial \cite{down:tomc}. Una vez descargado debemos configurar la variable de entorno \emph{CATALINA\_HOME} estableciendo la ruta de la carpeta de Tomcat. Además hay que incluir en el path la ruta de la carpeta \emph{bin} de Tomcat.

Terminada la configuración se puede realizar el proceso de su integración en Eclipse o desplegar la aplicación en la versión independiente.


\section{Compilación, instalación y ejecución del proyecto}

\subsection{Obtener código fuente}

Para obtener el código fuente desde GitHub se necesita clonar el repositorio.

Como se ha utilizado la herramienta de Git integrada en Eclipse se procede a explicar los pasos para realizar el proceso en este entorno.

\begin{enumerate}
\item Nos desplazamos a la vista de Git.
\imagen{GitView.png}{Vista de Git en Eclipse.}
\item Pulsamos el botón de clonar repositorio y pegamos la URL obtenida desde GitHub.
\imagen{GitClone1.png}{Paso 1 clonar repositorio.}
\item Pulsamos siguiente y seleccionamos la rama a clonar.
\imagen{GitClone2.png}{Paso 2 clonar repositorio.}
\item Por último seleccionamos el directorio donde queremos clonar el repositorio y pulsamos en \emph{Finish}.
\imagen{GitClone3.png}{Paso 3 clonar repositorio.}
\end{enumerate}

\subsection{Crear proyecto}

Una vez hemos clonado el repositorio en nuestro disco local debemos crear un proyecto para importar los ficheros al entorno de trabajo de Eclipse. Para ello creamos un proyecto en el directorio donde se ha clonado el repositorio anteriormente.

\imagen{CreateProject.png}{Creación de un proyecto en Eclipse.}

Una vez hemos creado el proyecto ya es posible editar o modificar todos los ficheros deseados teniendo en cuenta la estructura de un proyecto Maven \cite{maven:dir}.

\begin{itemize}
\item \textbf{src/main/java:} Código fuente de la aplicación.
\item \textbf{src/main/resources:} Recursos de la aplicación.
\item \textbf{src/main/webapp:} Código de la aplicación web.
\item \textbf{src/test/java:} Código de test.
\item \textbf{src/test/resources:} Recursos de los test.
\item \textbf{pom.xml:} Fichero para declarar dependencias de forma dinámica.
\end{itemize}


\subsection{Compilación y ejecución}

El proceso de compilación en Eclipse está configurado por defecto para realizarse automáticamente.

Antes de realizar la ejecución del proyecto es necesario incluir en el \emph{path} las dependencias de la carpeta \emph{lib}.

Para la ejecución del proyecto hay que seleccionar el proyecto y especificar la ejecución en servidor.

\imagen{RunAs.png}{Ejecutar en servidor.}

Seleccionamos Tomcat y automáticamente arranca el servidor y carga nuestra aplicación web.

Otra opción, si no se ha integrado el servidor Tomcat en Eclipse, es exportar la aplicación, copiar el archivo resultante en la carpeta específica del servidor y arrancarlo para que cargue nuestra aplicación.

Para realizar este proceso, se despliega el menú contextual del proyecto y se selecciona \emph{export}. En la siguiente pantalla se selecciona el formato en el que se exportará el proyecto, seleccionando WAR.

\imagen{Export1.png}{Paso 1 exportar proyecto.}

Al pulsar siguiente nos pide especificar la ruta donde colocar el fichero resultante. Se introduce y se finaliza la operación.

\imagen{Export2.png}{Paso 2 exportar proyecto.}

Una vez hemos obtenido el fichero, es necesario copiarlo en la ruta de nuestro servidor Tomcat.

./apache-tomcat-9.0.0.M17/webapps.

AL finalizar este paso, se procede a iniciar el servidor mediante el comando \emph{startup}.

\imagen{startup.png}{Arrancar el servidor.}

En el proceso de arranque, el servidor se encarga de construir la aplicación, quedando accesible desde la dirección local del servidor Tomcat.

\subsection{SonarQube}

SonarQube sirve para analizar la calidad del proyecto y ha sido utilizado en su versión local.

Para descargar SonarQube hay que dirigirse a su página oficial \cite{down:sonq}. Además hay que descargar Sonar Runner \cite{down:sonr}. Una vez que tenemos todo descargado hay que arrancar SonarQube dirigiendose a:

./sonarqube-6.4/bin/\{VersionSO\}

Y ejecutar StartSonar.bat en el caso de Windows.

\imagen{SonarQube.png}{SonarQube.}

Arrancado SonarQube podemos proceder a analizar el proyecto deseado. Para ello abrimos una terminal y nos situamos en el directorio del proyecto. Al tratarse de un proyecto Maven, ejecutamos los comandos \emph{mvn clean} para limpiar las versiones anteriores compiladas y \emph{mvn install} para compilar las clases.

\imagen{SonarRunner1.png}{Paso 1 ejecución SonarQube.}

Acabado este paso se puede proceder a analizar el proyecto mediante el comando \emph{mvn sonar:sonar}

\imagen{SonarRunner2.png}{Paso 2 ejecución SonarQube.}

Cuando termina el análisis tenemos disponible en la página local de SonarQube el resumen del análisis.

\imagen{SonarQubeAnalisys.png}{Resumen análisis SonarQube.}

El resumen completo puede verse en \href{https://sonarcloud.io/dashboard?id=UBUassistant}{Resumen SonarQube.}

\section{Pruebas del sistema}

Se han realizado dos tipos de pruebas para comprobar el correcto funcionamiento del sistema, pruebas unitarias de cada módulo y pruebas de integración e interfaz.

\subsection{Test unitarios}

Los test unitarios han sido realizados con JUnit 4, comprobando que cada módulo independiente del sistema tiene un funcionamiento correcto.

Se han realizado 21 test unitarios que realizan comprobaciones de los módulos de 7 clases que no necesitan un entorno web para funcionar. 

\imagen{TestUnitarios.png}{Test unitarios.}

Para su ejecución simplemente se seleccionan los test deseados y se ejecutan como JUnit.

\imagen{EjecucionTestUnitarios.png}{Ejecucion de una clase de test unitarios.}

\subsection{Test de interfaz}

Los test de interfaz han sido realizados mediante el framework de Selenium que otorga herramientas para realizar test de interfaz de forma automática sobre aplicaciones web. Estos test además comprueban la interacción de cada uno de los módulos testeados individualmente mediante los test unitarios.

Se ha realizado test que comprueban cada requisito funcional.

\imagen{TestInterfaz.png}{Test de interfaz.}



\tablaSinColores{Requisitos por test de interfaz.}
{L{4cm} L{8cm}}
{2}
{Requisitos por test de interfaz.}
{\textbf{Test} & \textbf{Requisito} \\}
{AdminAddCase.java 		&
 \begin{itemize}
 \item RF-7.1 Listar casos.
 \item RF-7.2 Añadir caso.
 \item RF-7.6 Ordenar casos.
 \end{itemize}  				
 \\
 AdminLearn.java 		&
 \begin{itemize}
 \item RF-8.1 Listar recomendaciones.
 \item RF-8.2 Aprender recomendación.
 \item RF-8.3 Descartar recomendación.
 \item RF-8.4 Exportar recomendaciones.
 \item RF-8.5 Ordenar recomendaciones.
 \end{itemize}  				
 \\
 AdminLog.java 		&
 \begin{itemize}
 \item RF-6.1 Listar log.
 \item RF-6.2 Limpiar log.
 \item RF-6.3 Exportar log.
 \item RF-6.4 Ordenar log.
 \end{itemize}  				
 \\
 AdminModifyCase.java 		&
 \begin{itemize}
 \item RF-7.1 Listar casos.
 \item RF-7.3 Editar caso.
 \item RF-7.4 Eliminar caso.
 \item RF-7.5 Exportar casos.
 \item RF-7.6 Ordenar casos.
 \end{itemize}  				
 \\
 EnterAdminCase.java 		&
 \begin{itemize}
 \item RF-6 Gestión del log.
 \item RF-7 Gestión de casos.
 \item RF-8 Gestión de aprendizaje.
 \end{itemize}  				
 \\
 MinimizeAndClose.java 		&
 \begin{itemize}
 \item RF-1 Reconocimiento de texto.
 \item RF-2 Búsqueda de respuesta.
 \item RF-3 Aprendizaje bajo supervisión.
 \item RF-4 Realizar recomendaciones.
 \item RF-5 Ofrecer mecanismo de valoración.
 \end{itemize}  				
 \\
 MultipleAnswer.java 		&
 \begin{itemize}
 \item RF-1 Reconocimiento de texto.
 \item RF-2 Búsqueda de respuesta.
 \item RF-4 Realizar recomendaciones.
 \item RF-5 Ofrecer mecanismo de valoración.
 \end{itemize}  				
 \\
 NoAnswer.java 		&
 \begin{itemize}
 	\item RF-1 Reconocimiento de texto.
 	\item RF-2 Búsqueda de respuesta.
 	\item RF-3 Aprendizaje bajo supervisión.
 	\item RF-4 Realizar recomendaciones.
 	\item RF-5 Ofrecer mecanismo de valoración.
 \end{itemize}  				
 \\
 ReservedAndSimpleAnswer.java 		&
 \begin{itemize}
 \item RF-1 Reconocimiento de texto.
 \item RF-2 Búsqueda de respuesta.
 \item RF-5 Ofrecer mecanismo de valoración.
 \end{itemize}  				
 \\
 SimpleAnswerAndStarBar.java 		&
 \begin{itemize}
 \item RF-1 Reconocimiento de texto.
 \item RF-2 Búsqueda de respuesta.
 \item RF-5 Ofrecer mecanismo de valoración.
 \end{itemize}  				
 \\}
\apendice{Documentación de usuario}

\section{Introducción}

En este apartado se van a detallar los requisitos necesarios para poder ejecutar la aplicación Android, así como la web de administración. A parte de esto, se especifica los pasos necesarios para la instalación y el funcionamiento de los mismos.

\section{Requisitos de usuarios}

\subsubsection{Aplicación Android}

Los requisitos para poder ejecutar la aplicación en Android son:

\begin{itemize}
	\tightlist
	\item 
	Dispositivo Android con versión 4.4 o superior.
	\item 
	Acceso a Internet (WiFi o datos móviles).
\end{itemize}

\subsubsection{Web de administración}

Los requisitos necesarios para poder utilizar la web de administración son:

\begin{itemize}
	\tightlist
	\item 
	Dispositivo con posibilidad de utilizar un navegador web.
	\item 
	Navegador web con JavaScript compatible y activado.
	\item 
	Acceso a Internet.
\end{itemize}

\section{Instalación}

\subsubsection{Aplicación Android}

Para poder instalar la aplicación APK, primero tenemos que activar la opción \textit{Orígenes desconocidos} en el dispositivo. Para ello accedemos a los ajustes del dispositivo, accedemos a la opción \textit{Seguridad} y activamos la opción \textit{Orígenes desconocidos}.

\imagenPeque{CapturaSeguridad.png}{Opciones de seguridad Android}

A continuación procedemos a la instalación de la aplicación, para ello primero tenemos que copiar el fichero APK a nuestro dispositivo. Una vez que termine la copia del fichero, desde nuestro dispositivo abrimos el fichero con un explorador de archivos. En la ventana que nos aparece, le damos a la opción \textit{Instalar}.

\imagenPeque{CapturaInstalacion.png}{Instalación UBUassistant}

Con estos pasos, la aplicación ya está instalada en el dispositivo Android.

\subsection{Web de administración}

Para poder utilizar la web de administración solamente es necesario tener instalado en el dispositivo en el que se vaya a utilizar un navegador web, ya que la aplicación de administración está alojado en un servidor.

Para acceder a la página de administración en el navegador hay que introducir la siguiente URL:

\begin{center}
	\textit{\textbf{http://IP\_SERVIDOR:8080/UBUassistant\_admin}}
\end{center}

\section{Manual del usuario}

En este manual se trata de guiar al usuario en el funcionamiento del asistente virtual con la aplicación Android, así como en la aplicación de administración web.

\subsection{Aplicación Android}

Lo primero que vemos al iniciar la aplicación es una ventana en la que aparecen el logo de la aplicación, junto al de la Universidad de Burgos. Lo siguiente que aparece es la interfaz del asistente virtual, la cual se asemeja a diferentes aplicaciones de mensajería para dispositivos móviles.

\imagenMediana{CapturaInterfaz.png}{Interfaz aplicación Android}

\newpage

En el caso de que los servidores de la aplicación no se encuentren disponibles, la aplicación se dirigrá automáticamente a una ventana en la que se muestra un mensaje de advertencia.

\imagenMediana{CapturaError.png}{Aviso de error de conexión}

En la parte inferior, encontramos una barra, la cual encontramos un cuadro donde podemos introducir el texto para interactuar con el asistente virtual. También nos encontramos con un botón para poder introducir el texto mediante voz y dos botones los cuales envían el texto al servidor.

\imagenMediana{CapturaVoz.png}{Reconocimiento de voz}

Dependiendo de la búsqueda que hagamos, podemos encontrar diferentes situaciones de respuesta.

\newpage

\subsubsection{Interacción con el asistente virtual}

Se puede interaccionar con el asistente virtual con mensajes como \textit{hola, buenos días, adiós, etc}.

\imagenMediana{CapturaSaludo.png}{Interacción con el asistente virtual}

\newpage

\subsubsection{Respuesta única}

Tras realizar una consulta con el asistente virtual, se pueden obtener diferentes situaciones de respuesta. Una de ellas, es que la búsqueda tenga una única respuesta. Un ejemplo de esto, puede ser la búsqueda del término \textit{correo}.

\imagenMediana{CapturaUnica.png}{Respuesta única}

\newpage

\subsubsection{Múltiples respuestas}

Otra de las posibles situaciones, es que la búsqueda que hemos realizado tenga múltiples respuestas. En esta situación, el asistente nos enseña todas las posibles respuestas que ha encontrado, para que se elija la que mas se ajusta al objetivo del usuario. Un ejemplo de este tipo de búsqueda es utilizando el término \textit{informática}.

\imagenMediana{CapturaMultiple.png}{Múltiples respuestas}

\newpage

\subsubsection{Sin respuesta}

La última posible situación que se puede llegar a dar, es que la búsqueda que se ha realizado no posee ninguna respuesta. En este caso, el asistente mostrará tres recomendaciones para ver si se ajustan al objetivo de búsqueda realizada por el usuario. Para obtener un ejemplo de este tipo de búsqueda, se puede utilizar el término \textit{Google}.

\imagenMediana{CapturaSin.png}{Pregunta sin respuesta}

\subsubsection{Enviar búsqueda y nueva búsqueda}

En la interfaz, podemos encontrar dos botones, uno con el texto \textit{ENVIAR} y otro con el texto \textit{ŃUEVO}.

El botón \textit{ENVIAR} sirve para enviar la respuesta siguiendo con las búsquedas que se han hecho con anterioridad, de esta manera, el usuario puede ir refinando la búsqueda hasta obtener la respuesta deseada.

\imagenMediana{CapturaEnviar.png}{Continuar búsqueda}

El botón \textit{NUEVO} sirve para reiniciar los términos de la búsqueda, de forma que el usuario puede empezar a realizar una búsqueda desde cero, no teniendo en cuenta el asistente los términos utilizados anteriormente.

\imagenMediana{CapturaNuevo.png}{Reiniciar búsqueda}

\subsubsection{Página web}

Cuando el usuario selecciona una respuesta otorgada por el asistente virtual, se nos abre la página web de la respuesta desde la propia aplicación, pudiendo interactuar con ella sin tener que salir de la aplicación.

\imagenMediana{CapturaWeb.png}{Respuesta consultada}

\subsubsection{Valoración}

Cuando el usuario termina de consultar la página web y vuelve al asistente, una ventana emergente, en la cual se le da la opción de valorar la respuesta obtenida.

\imagenMediana{CapturaValoracion.png}{Valoración de la respuesta}

\subsection{Aplicación web de administración}

Según accedemos a la aplicación de administración, lo que vemos es en la parte superior es el logo de UBUassistant. En la parte superior derecha, nos encontramos con un símbolo de un engranaje y llave inglesa, el cual nos dirige a la página de inicio de sesión de administrador. En la parte inferior izquierda, nos encontramos el logo de la aplicación, el cual si lo pulsamos, se abre el asistente, de tal forma que el administrador puede realizar todas las pruebas pertinentes sobre el funcionamiento del asistente.

\imagen{AdminInterfaz.png}{Interfaz de administración}

Cuando pulsamos sobre el logo situado en la parte derecha superior, se accede a la pantalla en la que el administrador tiene que introducir sus credenciales.

\imagen{AdminLogin.png}{Interfaz de inicio de sesión}

Una vez que el administrador ha iniciado sesión correctamente, automáticamente se abre la interfaz de administración. En la parte superior de la interfaz, se encuentran los botones para que el administrador pueda navegar hacia las diferentes vistas.

Como pantalla principal, aparece el \textit{log} del sistema. En esta ventana, el administrador puede visualizar todas las respuestas que se han abierto con el asistente en forma de tabla. El dato más importante que puede consultar el administrador, es la valoración que el usuario le ha dado a la respuesta.

\imagen{AdminLog.png}{Interfaz de \textit{log}}

En la interfaz de aprendizaje, podemos ver las recomendaciones de los usuarios para que el algoritmo aprenda en los casos de que no puede ofrecer respuestas. Estas mejoras vienen de cuando el algoritmo no encuentra respuesta y de las recomendaciones ofrecidas al usuario, éste abre una de ellas. A parte de ver las recomendaciones, el administrador desde esta página puede aceptar la recomendación, por lo que al asistente lo aprende en el momento; o desechar la recomendación de aprendizaje.

\imagen{AdminAprende.png}{Interfaz de aprendizaje}

En la opción \textit{Editor de casos}, el administrador se encuentra con dos opciones, \textit{Añadir caso} y \textit{Modificar caso}.

\imagenMediana{AdminCasosMenu.png}{Menú editor de casos}

\label{anadirCaso}
En el caso de \textit{Añadir caso}, la interfaz muestra un formulario, con el cual se puede añadir un nuevo caso a la base de datos. El formulario tiene unos campos que son necesarios para poder añadir un caso. Estos campos obligatorios son, \textit{Palabra Clave 1}, \textit{Categoría} y \textit{Respuesta}.

\imagen{AdminCasoAnadir.png}{Interfaz añadir caso}

En la opción \textit{Modificar casos}, nos aparece en forma de tabla todos los casos que posee el asistente virtual. Desde esta interfaz, el administrador puede editar o eliminar un caso.

\imagen{AdminCasoTabla.png}{Interfaz de administración de casos}

En el caso de que se elija la opción \textit{Editar}, se abre una nueva interfaz con un formulario igual al de \textit{Añadir caso} (\ref{anadirCaso}), teniendo que tener rellenados los mismos campos obligatorios. En este caso, el formulario trae campos rellenados según el caso que hayamos elegido para modificar.

\imagen{AdminCasoEditar.png}{Interfaz editar caso}

En el caso de que se elija la opción \textit{Eliminar}, aparece un mensaje de aviso, para preguntarnos si estamos seguros de querer borrar el caso.

\imagen{AdminCasoEliminar.png}{Aviso de borrado de caso}

A parte de todas estas vistas que se han comentado, en el caso de que los datos se muestran en forma de tabla, el administrador tiene la opción de extraer la tabla en diferentes formatos, Esto se puede realizar pulsando el botón \textit{Guardar como}, situado en la parte superior izquierda de las tablas.

\imagen{AdminExport.png}{Exportar tabla}

Otra función importante de las tablas es que son ordenables. De esta forma, el administrador, puede ver los datos de la forma que crea más conveniente. Para ordenar la tabla, solamente hay que pulsar sobre la cabecera de la columna que se quiere ordenar. La primera vez, la tabla se ordena de forma ascendente o de A-Z. La segunda vez, la tabla se ordena de forma descendente o de Z-A.


\bibliographystyle{IEEEtran}
\bibliography{IEEEabrv,bibliografiaAnexos}

\end{document}
